\chapter{Introducción}

\section{Sobre este manual}

Este manual está dirigido a nuevos desarrolladores que quieren trabajar en este proyecto. Su objetivo es familiarizar a los nuevos desarrolladores con el proyecto, y brindar tanto las herramientas como la información necesaria para que puedan desarrollar sobre el proyecto de manera rápida y efectiva.

Este manual contiene el contexto del proyecto y los objetivos que quiere alcanzar (Capítulo 1), así como las funcionalidades necesarias para alcanzar dichos objetivos (Capítulo 2). También, se incluye la información sobre cómo se diseño el sistema, y cómo están implementadas las funcionalidades requeridas (Capítulo 3). Además, contiene guías para configurar el entorno de desarrollo y para desplegar el sistema (Capítulo 4). Finalmente, se incluyen reflexiones útiles sobre el proceso de desarrollo, las limitaciones del proyecto y las oportunidades de mejora que existen (Capítulo 5).

\section{Descripción y Alcance del proyecto}
\subsection{Antecedentes}

Este proyecto surge como respuesta a una necesidad directa expresada por el cliente, quien en la actualidad usa distintos lectores de PDF y manifiesta una insatisfacción con estas herramientas debido a sus interfaces sobrecargadas de utilidades, barras de herramientas y menús poco intuitivos. Además, también menciona experiencias negativas con lectores alternativos, que son más livianos, pero tienen algunos problemas de estabilidad, entre otras cosas.

Como referencia de estilo visual el cliente resaltó la interfaz que utiliza un lector desarrollado por Augmented Text Company (\url{https://www.augmentedtext.info/reader}), que corresponde a un visor sin barra de título, sin bordes intrusivos y con un estilo minimalista blanco sobre blanco agradable a la vista.

\subsubsection{Beneficiario}
El beneficiario directo del proyecto es el cliente Rodolfo Mora, quien en su labor académica maneja de forma constante documentos en formato PDF (lecturas, presentaciones exportadas y documentos legales que requieren firma), y que actualmente no cuenta con una herramienta
que combine una interfaz minimalista con las funcionalidades específicas que necesita.

\subsubsection{Importancia}
La importancia del proyecto radica en ofrecer una alternativa amigable a las aplicaciones de lectura y edición de PDF existentes en el mercado que, según lo expresado por el cliente, priorizan la cantidad de funcionalidades visibles por encima de la experiencia de lectura sobrecargando la interfaz y dificultando la concentración en el contenido. 

\subsubsection{Grupos o comunidades beneficiadas}

El principal beneficiado es el cliente. Sin embargo, dado que se contempla como un proyecto de software libre, cualquier usuario de Windows (y eventualmente Linux) que busque un lector y editor de PDF minimalista podrá beneficiarse de la herramienta desarrollada.

\subsubsection{Necesidades solventadas}

Interfaces sobrecargadas y poco intuitivas de los lectores de PDF actualmente disponibles (Adobe Acrobat, Foxit, entre otros), que dificultan una experiencia de lectura centrada en el contenido.
	

\subsection{Objetivos}

%
% Los objetivos del proyecto deben englobar el alcance total. Todo objetivo debe estar redactado respetando la siguiente estructura:

% Verbo infinitivo: Representa una acción que el equipo de desarrollo ejecutará durante el proyecto. Debe ser una acción concreta y relevante para completar el proyecto

% Objeto: Indica el objeto sobre el cual se aplicará la acción. Típicamente representa el producto esperado de la acción

% Modo: Indica parámetros por los cuales se medirá la calidad y completitud del resultado. 

% Ejemplo: Desarrollar un sistema de facturación electrónica que provea una interfaz amigable para el usuario y cumpla con los estándares impuestos por el Ministerio de Hacienda de Costa Rica. 

% Sólo puede existir un verbo infinitivo en el objetivo
% Debe ser posible medir el nivel de cumplimiento del objetivo, es decir, debe ser posible decir en cualquier momento del proyecto, cual porcentaje del objetivo se ha cumplido.

\begin{itemize}
    \item \textbf{Objetivo General}
    \begin{itemize}
        \item Desarrollar una aplicación de escritorio para Windows que
        permita visualizar y modificar documentos PDF mediante una interfaz minimalista, limpia y clara.
    \end{itemize}
    \item \textbf{Objetivos Específicos}
    \begin{itemize}
        \item Construir un visualizador de documentos pdf que muestre las páginas del documento al usuario.
        \item Desarrollar funcionalidades para hacer modificaciones simples al pdf, como eliminar, compaginar y mezclar páginas, y que se guarden en el mismo archivo o uno nuevo.
        \item Implementar la búsqueda, selección y copiado de texto en un pdf que se está visualizando.
        \item Desarrollar un modo presentación para el pdf que se está visualizando.
        \item Diseñar una interfaz minimalista con color blanco que muestre las funcionalidades de la aplicación solo cuando sea necesario, intentando mantener poca información en la pantalla para no saturar al usuario.
    \end{itemize}
\end{itemize}

\subsection{Contacto}
Para consultas y soporte, contacte a: samuelzv18@gmail.com

